Modern software development increasingly relies on complex microservice architectures, often utilizing robust frameworks such as Spring Boot and event-streaming platforms like Apache Kafka. While these technologies accelerate functional development, they often introduce complexity that makes rigorous quality assurance and formal verification difficult. Open-source repositories frequently demonstrate functional correctness but lack the software dependability layers—such as mutation testing, strict static analysis, and formal specifications—required for production-grade systems.

This paper presents a case study on retrofitting dependability measures onto an existing open-source artifact. We selected the \repourl\ repository \cite{ecommerceRepo}, a representative microservices application implementing an e-commerce backend. While the original repository provided a functional implementation of shopping cart and order processing logic, it lacked a comprehensive automated testing strategy and formal behavioral specifications.

Our work follows a systematic approach to software hardening, addressing the following key areas:

\begin{enumerate}
    \item \textbf{Environmental Stabilization:} We addressed immediate build reproducibility issues by upgrading the Docker base image to \texttt{eclipse-temurin:17-jdk} and pinning Kafka/Zookeeper versions to resolve startup failures in the containerized environment.
    \item \textbf{Automated Quality Gates:} We established a continuous integration (CI) pipeline using GitHub Actions to enforce build verification on every push.
    \item \textbf{Static and Dynamic Analysis:} We integrated a suite of static analysis tools (SpotBugs, PMD, Checkstyle) and dynamic coverage analysis (JaCoCo) to establish a quality baseline.
    \item \textbf{Advanced Verification:} We extended the test suite based on coverage data and applied Mutation Testing (PiTest) to assess the robustness of our tests. Additionally, we applied Java Modeling Language (JML) annotations to core services to formalize system behavior, identifying architectural constraints imposed by framework-heavy designs (e.g., Lombok) that complicate formal verification.
\end{enumerate}

The remainder of this paper details the methodology used to integrate these tools, presents the empirical metrics gathered (coverage improvements, mutation scores), and discusses the challenges of applying formal specifications to modern Spring Boot applications.